\documentclass[11pt,a4paper]{article}
\usepackage{hyperref}
\usepackage{listings}
\usepackage{xcolor}

\title{Report PCD 2025-2026 \\ Distributed Critical Sections with a Message-Oriented Middleware (Assignment 4 - Exercise 3)}
\author{Francesco}
\date{\today}

\lstset{
  basicstyle=\ttfamily\small,
  keywordstyle=\color{blue},
  stringstyle=\color{red},
  commentstyle=\color{green!60!black},
  numbers=left,
  numberstyle=\tiny,
  frame=single,
  breaklines=true
}

\begin{document}

\maketitle

\section{Introduction}
This project implements a high-level middleware for distributed critical sections using Java and RabbitMQ.
The goal is to let independent processes coordinate access to a shared critical section without knowing
anything about each other beyond the critical-section name.

\section{Analysis}
The problem is a distributed mutual exclusion problem with crash recovery constraints:
\begin{itemize}
    \item \textbf{Mutual exclusion}: at most one process may enter the critical section at a time.
    \item \textbf{Crash recovery}: if the holder crashes while inside the critical section, the token must be
    returned automatically by the broker.
    \item \textbf{Bootstrap safety}: initialization must never leave the system in a state where the critical
    section is marked initialized but the token does not exist.
    \item \textbf{Independence}: processes must not coordinate through local static state, shared files, or any
    direct inter-process channel.
    \item \textbf{Broker semantics}: the solution must rely on RabbitMQ message passing and broker-provided
    primitives only.
\end{itemize}

\section{Design}
The middleware is implemented as a token queue protocol:
\begin{itemize}
    \item \textbf{Token queue}: each critical section is mapped to a durable queue named
    \texttt{cs\_token\_<csName>}.
    \item \textbf{Token message}: the queue contains exactly one persistent token when the critical section is
    available.
    \item \textbf{Acquire}: \texttt{enter()} creates a consumer with manual acknowledgement and blocks until the
    token delivery arrives.
    \item \textbf{Release}: \texttt{exit()} requeues the same delivery using \texttt{basicNack(..., requeue = true)}
    after the local consumer has been canceled.
    \item \textbf{Crash-safe bootstrap}: a temporary exclusive queue \texttt{cs\_bootstrap\_lock\_<csName>} serializes
    initialization. The broker is queried passively, and the token is published only when the queue is empty and
    has no consumers.
\end{itemize}

This design keeps the token itself as the only durable state that matters. The bootstrap lock queue exists only
while a process is deciding whether the token must be created, so no initialization marker can outlive the token.
If a process crashes while holding the token, RabbitMQ requeues the unacknowledged delivery when the connection
or channel closes.

The implementation also preserves lifecycle clarity:
\begin{itemize}
    \item the middleware may be created from an existing connection or by creating its own connection;
    \item \texttt{close()} closes the channel and, when appropriate, the connection;
    \item the public API is intentionally small: \texttt{enter()}, \texttt{exit()}, and \texttt{close()}.
\end{itemize}

\section{Validation}
Validation is covered by JUnit tests that exercise the full protocol:
\begin{enumerate}
    \item concurrent initialization by multiple middleware instances;
    \item exactly one token created for a critical section;
    \item interrupted bootstrap and later recovery by another process;
    \item crash while holding the token and automatic requeueing by RabbitMQ;
    \item mutual exclusion with multiple concurrent processes;
    \item repeated acquire/release cycles;
    \item rejection of release without ownership and of double release;
    \item independence of different critical-section names;
    \item shutdown cleanup without losing the token.
\end{enumerate}

\section{Conclusions}
The resulting middleware provides a simple and robust distributed critical-section service built entirely on
RabbitMQ message passing. The token queue acts as the single source of truth, the bootstrap lock queue prevents
duplicate initialization, and manual acknowledgements ensure the token returns automatically after failures.
This keeps the solution easy to reason about while still handling the main distributed-systems concerns of
coordination, crash recovery, and resource cleanup.

\end{document}
